%%%% default header %%%%
%REQUIRES packages: amssymb, pifont
\usepackage{xspace}
%%custom marcros

% common stuff
%\def\ie{{\em i.e.~}}
%\def\eg{{\em e.g.~}}
\def\ie{i.e.~}
\def\eg{e.g.~}
\def\st{s{.}t~}
\def\cf{{\em c.f.~}}
\def\etc{{\em etc.}}
\def\aka{{a.k.a~}}
\def\wrt{{\em w.r.t~}}
\def\resp{{\em resp.}}
\def\esp{{\em esp.~}}

% % abbreviations
\def\abbr#1{{\em abbr.} #1}
\def\abbrv#1#2{\textbf{#1} (\textbf{#2})}
% USED for textsc-typed abbreviation: similar to \abbrv but does not bolden #2
\def\abbrvs#1#2{\noindent{\textbf{#1} (#2)}}

% % logic
\def\ifff{\noindent\ensuremath{\leftrightarrow}}

% general function for names, e.g. \Name{JDomainApp}
\def\Name#1{{\footnotesize \textsc{#1}}}
% same as \Name but using standard font
\def\name#1{\textsc{#1}}

% annotation symbol, <<a>>
\def\anosymbol{\noindent{\footnotesize \it \ensuremath{\left\langle\!\left\langle{a}\right\rangle\!\right\rangle}}}

% complexity notation
\def\bigO{\textit{O}}
%string literals
\def\str#1{{\tt `#1'}}
\def\attribval#1#2#3{\noindent\ensuremath{\tt #1{.}#2{=}#3}}
% double quoted string in math mode
\def\strq#1{\noindent\ensuremath{``\tt #1"}}

% for writing Java code in-line with text
\newcommand\jcode[1]{\noindent\ensuremath{{\tt #1}}} % \texttt{#1}
% mutiple line set definition: {x | P}
% #1: x, #2: P, #3: width
\newcommand\Set[3]{\ensuremath{\{\text{#1 $|$ \parbox[t]{#3}{#2\}}}}}
% set with code-specific font
\def\set#1{\noindent\ensuremath{\left\lbrace\code{#1}\right\rbrace}}
%set with standard math font
\def\sets#1{\noindent\ensuremath{\left\lbrace#1\right\rbrace}}
% power set of a set
\def\powerset#1{\noindent\ensuremath{{\mathcal P}(#1)}}
% math tuple with mathsf font, e.g. <x,y> (mathsf)
\def\tuple#1{\noindent\ensuremath{\left\langle\code{#1}\right\rangle}}
% tuple with normal font, e.g. <a,b> (normal)
\def\tuplex#1{\noindent\ensuremath{\left\langle#1\right\rangle}}
% tuple with type prefix (normal font), e.g: A:<a,b> (normal font)
% 1: tuple type (e.g. A), type separator (e.g ':'), tuple elements (e.g. a,b)
\newcommand\tupleT[3]{\noindent\ensuremath{#1\!{#2}\!\left\langle#3\right\rangle}}
% variable (in listing-typed font) with subscripts, e.g. v_1, etc.
\newcommand\var[2]{\noindent\lstinline{#1}$_#2$}
% % math box (used inside tables)
%1: width, 2: content, 3: vertical ident at top
\def\rulebox#1#2#3{ %
\rule{0pt}{#3} %
\begin{minipage}[t]{#1}\raggedright %
\noindent#2\end{minipage} %
}

% % % FUNCTIONS 
% function definition: #1: func. name, #2: domain, #3: range
\newcommand\funcdef[3]{\noindent\ensuremath{{\tt #1}{:~}#2 \rightarrow #3}}
%
%hdm stuff %
%\def\hdmopenscheme{\mbox{\ensuremath{\langle\!\langle}}}
%\def\hdmclosescheme{\mbox{\ensuremath{\rangle\!\rangle}}}
%% denotes an HDM scheme, use in maths mode
%\def\hdm#1{\hdmopenscheme#1\hdmclosescheme}
%\def\hdmexl{\ensuremath{\mathbin{\not\!\cap}}}
%\def\hdminc{\ensuremath{\mathbin{\subseteq}}}
%\def\hdmuni{\ensuremath{\mathbin{\cup}}}
%END hdm stuff %
%
\def\courseman{\textsc{CourseMan}\xspace}
\def\jdomainapp{\textsc{jDomainApp}\xspace}
\def\orderman{\textsc{OrderMan}\xspace}
\def\processman{\textsc{ProcessMan}\xspace}
\def\domainapptool{{\textsc{DomainAppTool}}\xspace}
\def\dcsl{\textsc{DCSL}\xspace}
\def\mccl{\textsc{MCCL}\xspace}
\def\agcl{\textsc{AGCL}\xspace}
\def\agl{\textsc{AGL}\xspace}
\def\agldcsl{\textsc{AGL+}\xspace}
%
%\def\code#1{\noindent\ensuremath{\mathsf{#1}}}
% % % bold and italic
\def\bi#1{\textbf{\textit{#1}}}
\def\img#1{\imgscale{#1}{0.7}}
\def\imgscale#1#2{\includegraphics[scale=#2]{images/#1.eps}}
%
% indent by #1
%\def\indent#1{\noindent\parbox{#1}{~}}
% tab 
%\def\tab{\hspace{2em}}
\def\parhead#1{\noindent\textbf{#1}} % paragraph heading
% paragraph heading using mini page (with next line supported)
% 1: text width, 2: content 
\def\parheadbox#1#2{\noindent\begin{minipage}{#1}~\\[0.2em]\textbf{#2}\\[-0.5em]\end{minipage}} 
% similar to \parheadbox except that it uses italic rather than bold font
% 1: text width, 2: content 
\def\parheadboxit#1#2{\noindent\begin{minipage}{#1}~\\[0.2em]\textit{#2}\\[-0.5em]
\end{minipage}} 

% bold and italic
\newcommand\bfit[1]{\textbf{\textit{#1}}}

% % Table stuff
% column header with custom alignment
% #1: column align (e.g. c|), #2 column text 
\def\colheader#1#2{\multicolumn{1}{#1}{#2}}
% wrapped column header with custom alignment
% #1: column align spec (e.g. c|) 
% #2: header width (e.g. 1cm), #3 column text (possibly with nextline character(s))
\def\colheaderwrap#1#2#3{\multicolumn{1}{#1}{\parbox{#2}{\vspace{0.2em} #3 \\[0.2em]}}}
% wrapped table cell text with custom alignment
% #1: align spec (e.g. c|) 
% #2: cell width (e.g. 1cm), #3 cell text (possibly with nextline character(s))
\def\cellwrap#1#2#3{\multicolumn{1}{#1}{\parbox{#2}{\vspace{0.2em} #3 \\[0.2em]}}}

% % END: table stuff

% rotate text box
% #1: text, #2: rotation angle
\def\rotatetext#1#2{\noindent\rotatebox[origin=c]{#2}{#1}}

% reflect: model reflect domain class
\def\reflect{\curvearrowleft}

% % special symbol tables 
\def\Bigdiamond{\lozenge}
\def\Bigfilleddiamond{\blacklozenge}

% % END symbol tables

% % circled numbers 
% % Requires: package pifont
\def\circledOne{\ding{182}}
\def\circledTwo{\ding{183}}
\def\circledThree{\ding{184}}
\def\circledFour{\ding{185}}
\def\circledFive{\ding{186}}
\def\circledSix{\ding{187}}
\def\circledSeven{\ding{188}}
\def\circledEight{\ding{189}}
\def\circledNine{\ding{190}}
\def\circledTen{\ding{191}}
% requires package tikz
\def\circleHalfFull{
\begin{tikzpicture}
\draw (0,0) arc (0:180:0.1cm); % top half (an arc)
\filldraw (0,0) arc (0:-180:0.1cm); % bottom half (a filled arc)
\end{tikzpicture}
}
\def\circleHalfFullCell{
\begin{tikzpicture}
\draw (0,0) arc (0:180:0.1cm); % top half (an arc)
\filldraw (0,0) arc (0:-180:0.1cm); % bottom half (a filled arc)
\end{tikzpicture}\vspace{-2pt}
}
% unfilled numbers
\def\unfilledCircledOne{\ding{172}} % 192,...
\def\unfilledCircledTwo{\ding{173}}
\def\unfilledCircledThree{\ding{174}}
\def\unfilledCircledFour{\ding{175}}
\def\unfilledCircledFive{\ding{176}}
\def\unfilledCircledSix{\ding{177}}
\def\unfilledCircledSeven{\ding{178}}
\def\unfilledCircledEight{\ding{179}}
\def\unfilledCircledNine{\ding{180}}
\def\unfilledCircledTen{\ding{181}}
% filled numbers
\def\fone{\ding{202}} % 182,...
\def\ftwo{\ding{203}}
\def\fthree{\ding{204}}
\def\ffour{\ding{205}}
\def\ffive{\ding{206}}
\def\fsix{\ding{207}}
\def\fseven{\ding{208}}
\def\feight{\ding{209}}
\def\fnine{\ding{210}}
\def\ften{\ding{211}}
%
\def\tick{{\normalsize \ding{51}}}
\def\tickbold{{\normalsize \ding{52}}}
\def\cross{{\small \ding{56}}}
\def\tickcell{{\normalsize \ding{51}}\vspace{-2pt}}
\def\tickboldcell{{\normalsize \ding{52}}\vspace{-2pt}}
\def\crosscell{{\small \ding{56}}\vspace{-3pt}}
% end requires pifont

%caligraphy
\def\calA{\noindent\ensuremath{\mathcal{A}}}
\def\calN{\noindent\ensuremath{\mathcal{N}}}
\def\calT{\noindent\ensuremath{\mathcal{T}}}
\def\calR{\noindent\ensuremath{\mathcal{R}}}

%special symbols
\def\Tau{\noindent\ensuremath{\tau}}
% language extension symbol
\def\lext{\noindent\rotatebox[origin=c]{270}{\ensuremath{\Lsh}}}

% % number sets: requires either amsfonts or amssymb
% natural numbers
\def\Nat{\noindent\ensuremath{\mathbb{N}}}
% integer numbers
\def\Int{\noindent\ensuremath{\mathbb{Z}}}
% rational numbers
\def\Rat{\noindent\ensuremath{\mathbb{Q}}}
% irrational numbers
\def\Irat{\noindent\ensuremath{\mathbb{I}}}
% real numbers
\def\Real{\noindent\ensuremath{\mathbb{R}}}
% complex numbers
\def\Complex{\noindent\ensuremath{\mathbb{C}}}
% prime numbers
\def\Prime{\noindent\ensuremath{\mathbb{P}}}